\documentclass[11pt, letterpaper]{article}
\usepackage[utf8]{inputenc}
\usepackage[T1]{fontenc}
\usepackage{amsmath}
\usepackage{amsfonts}
\usepackage{amssymb}
\usepackage{geometry}
\usepackage{verbatim}

\geometry{letterpaper, margin=1in}

\title{\textbf{Problem C: Koyomi's Revisions}}
\author{Time Limit: 3.0s | Memory Limit: 512MB}
\date{}

\begin{document}

\maketitle

\section*{Description}

"I wrote this... but does it actually mean anything new?"

Koyomi Kanou is struggling with the script for the Student Council play. She has written a long draft, represented as a string $S$ of lowercase English letters. She is constantly analyzing specific sections of the script to ensure they are rich in content.

Koyomi defines the **Originality** of a section $S[L \dots R]$ (the substring starting at index $L$ and ending at $R$) as the number of **distinct substrings** that appear wholly within that range.

You must answer $Q$ queries. For each query $(L, R)$, calculate the number of distinct strings $T$ such that $T$ occurs as a substring of $S[L \dots R]$.

\section*{Input}

The first line contains a string $S$ of length $N$ ($1 \le N \le 2 \cdot 10^5$).
The second line contains an integer $Q$ ($1 \le Q \le 2 \cdot 10^5$).
The next $Q$ lines each contain two integers $L$ and $R$ ($1 \le L \le R \le N$).

\section*{Output}

For each query, output the Originality (count of distinct substrings) of the range $S[L \dots R]$.

\section*{Examples}

\subsection*{Example 1}
\textbf{Input:}
\begin{verbatim}
abacaba
3
1 3
2 4
1 7
\end{verbatim}

\textbf{Output:}
\begin{verbatim}
5
6
21
\end{verbatim}

\textbf{Explanation:}
Query 1 ($S[1 \dots 3]$ = "aba"):
Distinct substrings: "a", "b", "ab", "ba", "aba". (Count: 5).

Query 2 ($S[2 \dots 4]$ = "bac"):
Distinct substrings: "b", "a", "c", "ba", "ac", "bac". (Wait, "bac" is length 3. Substrings: b, a, c, ba, ac, bac. Count 6? Let's check).
"b", "a", "c", "ba", "ac", "bac". Yes 6.
Wait, output says 5?
Ah, example output in my head was wrong. Let's trace carefully.
"bac": "b", "a", "c", "ba", "ac", "bac". Total 6.
Why did I write 5? Maybe I missed one.
Let's trust the solution code to generate correct outputs.

Query 3 ($S[1 \dots 7]$ = "abacaba"):
This is the full string.
Total distinct substrings of "abacaba" is 21.

\end{document}
